% =============================================================================
\section{Discussion}
\label{sec:discussion}
% =============================================================================

\subsection{Problems met and how they were solved}

\paragraph{Button bouncing.} A mechanical contact bounces for several milliseconds, so a raw read
would turn one press into several RUN/PAUSE toggles. The integrating debouncer
(Section~\ref{sec:impl}) only accepts a new level after 30\,ms of stability and emits exactly one
latched edge. On the K1 key, the board's RC filter (10\,k$\Omega$/100\,nF, $\tau = 1$\,ms) helps, but the
software filter is what makes the behaviour deterministic (HAL-01, TC20).

\paragraph{Timer handling.} The first idea of many students is ``PAUSE stops the timer''. The
specification says the opposite. Making the timer a context variable decremented by the \emph{same}
handler in RUNNING and PAUSED removed the risk of the two paths diverging. A second subtle point is the
end condition: the transition is taken on the tick that goes from 1 to 0, so the cycle lasts exactly
$T$ ticks and the FSM never sits in RUNNING with $t = 0$.

\paragraph{Avoiding blocking delays.} With \code{HAL_Delay()} the LED blinking alone would block the CPU
half of the time and inputs would be missed. Every delay was replaced by a counter advanced by the 1\,ms
tick. The catch-up loop (\code{while (processed != now)}) keeps the timing exact even when the main loop
is temporarily slow, e.g.\ while a full LCD frame is sent over SPI.

\paragraph{ISR / main-loop race.} An earlier version updated the FSM from inside the SysTick ISR while the
main loop also dispatched button events, so two contexts could modify the FSM at the same time. The final
design keeps the ISR trivial (\code{g_millis++}) and runs \emph{all} FSM work in the main loop, which
removes the race without disabling interrupts.

\paragraph{Double-press STOP.} The difficulty is to separate ``two presses'' from ``two unrelated
presses''. A 1.5\,s sliding window implemented with a countdown solves it, and the boundary is tested at
1499\,ms and 1501\,ms. On the board, a single key also had to produce two STOPs; the two-stage hold
(0.6\,s, 1.2\,s) does this and the second stage is guaranteed to fall inside the window.

\paragraph{One key for seven inputs.} The WeAct board exposes only K1 to software (BOOT0 is not a GPIO on
the H7). A gesture decoder (click, double-click, two-stage hold) plus state-dependent mapping gives every
command, without changing the FSM. Its cost is a 300\,ms delay before a single click is confirmed
(to distinguish it from a double-click).

\paragraph{State transitions.} Writing the transition table first (Table~\ref{tab:transitions}) and then
one \code{enter_xxx()} function per state made the LED rules impossible to forget on any incoming
transition. Test TC-27 enumerates all (state, event) pairs to check that unspecified events are ignored.

\paragraph{LCD refresh.} Redrawing 25.6\,KB at every loop iteration would waste time and cause flicker.
The firmware builds a small ``view'' structure and redraws only if it differs from the last drawn one, so
the panel is refreshed at most a few times per second. The backlight faded out when driven as a plain GPIO
on this board; driving PE10 with a 10\,kHz PWM from TIM1 fixed it.

\subsection{Limitations}

\begin{itemize}
  \item \textbf{No real coin acceptor.} Coins are simulated with K1 gestures, external buttons or UART keys.
        The pulse-train decoder for a real acceptor is implemented and unit-tested but not wired.
  \item \textbf{No real motor or relays.} Motor, pump, valve and door-lock commands are only displayed. A real
        machine would need relay/driver hardware and feedback sensors.
  \item \textbf{Only one physical LED.} RLED and BLED are drawn on the LCD; the PE3 LED shows their OR.
        Two separate LEDs could be added on free GPIOs without software changes (only the callbacks).
  \item \textbf{Error state is simulated.} Faults are injected by a K1 double-click; there is no real door
        switch, water-level or current sensor.
  \item \textbf{Time scaling for the demo.} The default board build uses a 30\,s cycle so the whole cycle can be
        shown; the 30-minute behaviour is verified by the host tests and by the \code{FULL_CYCLE=1} build.
  \item \textbf{No power-loss persistence.} A reset during a paid cycle loses the remaining time; a real
        product would store it in backup SRAM or flash.
  \item \textbf{Front-panel extensions.} The coin staging/confirm step and the ``change'' notice on the board
        go beyond the statement (which says READY as soon as 50\cent{} is reached and no refund). They are a
        UI layer only; the FSM follows the statement and its tests are unchanged.
\end{itemize}
